\sscp{Minor Austronesian incursion: from Sulawesi into central Maluku}{15-06-02}
In \crf{ch:Tal}, we provided various kinds of evidence that Taliabu,
Soboyo, and Kadai spoken in the Sula archipelago in western Maluku
near Sulawesi share a number of features in their historical
phonologies with Banggai and other \tsc{Celebic} languages to the west that are not shared
with any other \tsc{Sula-Buru},
\tsc{Ambon-Seram} or \tsc{Seram-Tanimbar-Bomberai} languages.
Particularly strong evidence is seen in sharing PMP *ə > \it{o},
*-ay > \it{e}, *-aw > \it{o} with \tsc{Celebic} languages,
but not with any other Austronesian languages of geographical central Maluku.

It was also shown that the \tsc{Taliabu} languages have had a
strong influence on the morphology
and lexicon of the nearby \tsc{Sula-Buru} languages of Mangole
and Sanana, a reduced influence on Lisela,
and very little on Buru and Ambelau.

\citet[130]{Bloyd2020} observes, “Sula was a casual swim’s distance
from Sulawesi during the height of the Ice Age.” \frf{02-04}
shows the shallow seas connecting Sulawesi
and Taliabu via Banggai.
In \srf{14-01} it was noted that the \it{Strigocuscus pelengensis}
‘Banggai (dwarf) cuscus’, is also found throughout the Sula archipelago,
but not across the gap to Buru.

This incursion runs counter to \citeauthor{Blust2012b}'s (\citeyear[274]{Blust2012b}) claim
that \citeauthor{Schapper2011}’s (\citeyear{Schapper2011}) scenario about marsupial terms,
“[\ldots] implies an entry via Sulawesi into the Lesser Sundas,
with subsequent eastward movement into the Moluccas.
If this were the case,
we would expect some linguistic trace in Sulawesi,
but there is none; despite the speculations in \citet{DonohueGrimes2008}
about possible subgrouping ties between at least some of the
languages of Sulawesi and languages in eastern Indonesia,
closer scrutiny of the evidence reveals nothing of substance
to support this idea.”

There are a bundle of phonological features that tie the languages
of \tsc{Taliabu} to Sulawesi languages, rather than to other languages in Maluku.
There are also a bundle of morphological features
and the very different lexicons that separate
the \tsc{Taliabu} languages from \tsc{Sula-Buru}.
It was a relatively minor incursion that does not appear to have
penetrated very far into Maluku.